% Sinh từ cùng nguồn nội dung với bản Word.
\section{Introduction}

\subsection{Background and Motivation}

Trong lắp ráp, tổng thời gian hoàn thành một chu kỳ phản ánh năng suất chung nhưng không cho biết thời gian được phân bổ thế nào giữa các bước như lấy linh kiện, định vị, lắp ghép hay dùng dụng cụ. Muốn biết vì sao một chu kỳ kéo dài, cần nhận ra từng thao tác, mốc bắt đầu, kết thúc của nó và các khoảng chờ hay thao tác lặp.

Hiện nay việc này thường làm bằng cách quan sát hoặc xem lại video, bấm giờ và ghi chép. Cách làm tốn thời gian khi số video tăng, và mỗi người đánh giá có thể chọn mốc bắt đầu, kết thúc khác nhau nên số liệu kém nhất quán. Một công cụ tự động nhận dạng và đo thời gian từng bước giúp giảm việc thủ công và tạo cơ sở so sánh thống nhất.

Ước lượng tư thế người cho phép biểu diễn chuyển động bằng chuỗi điểm khớp, từ đó mô hình học có thể nhận dạng thao tác và ranh giới thời gian của chúng. Tuy nhiên, nhận dạng thao tác lắp ráp có những khó khăn riêng: nhiều bước có động tác gần giống nhau nhưng khác linh kiện, một số thao tác lặp lại liên tiếp trong cùng chu kỳ, và bàn tay thường bị che khuất. Sau khi đo được thời gian, kết quả còn cần được đặt trong ngữ cảnh của quy trình chuẩn — vốn nằm trong tài liệu hướng dẫn công việc và SOP — để người phụ trách biết bước nào lệch chuẩn và nên cải tiến theo hướng nào. Dự án vì vậy kết hợp phân đoạn hành động từ chuỗi khung xương với truy xuất tăng cường sinh (RAG) để tạo báo cáo có căn cứ và truy vết được, phục vụ hoạt động Kaizen.

\subsection{Problem Statement}

Đầu vào của hệ thống là video RGB ghi một chu kỳ lắp ráp bộ bánh răng hành tinh tại một trạm cố định, định nghĩa 12 bước thao tác và tập tài liệu tham chiếu. Nhiệm vụ thứ nhất là học có giám sát: dự đoán nhãn bước và lần lặp tại từng khung hình từ chuỗi điểm khớp, rồi chuyển thành các đoạn có mốc bắt đầu, kết thúc. Nhiệm vụ thứ hai là đánh giá: tính thời lượng từng bước, đối chiếu với thời gian chuẩn truy xuất từ tài liệu, phát hiện các bước lệch ngưỡng và sinh khuyến nghị cải tiến có trích dẫn. Mọi phép tính thời gian và chênh lệch do chương trình thực hiện; mô hình ngôn ngữ chỉ diễn giải các số liệu đã tính và đoạn tài liệu đã truy xuất.

Câu hỏi nghiên cứu trọng tâm là: (i) mô hình phân đoạn học từ khung xương 2D có khái quát sang người thực hiện chưa xuất hiện khi huấn luyện hay không, và những yếu tố nào của dữ liệu giới hạn khả năng đó; (ii) sai số phân đoạn ảnh hưởng thế nào đến các phát hiện và khuyến nghị mà tầng RAG tạo ra.

\subsection{Objectives and Scope}

Mục tiêu là xây dựng và đánh giá thực nghiệm một luồng xử lý thống nhất từ video đến báo cáo khuyến nghị, gồm: trích xuất điểm khớp, phân đoạn bước bằng MS-TCN, đo thời lượng, truy xuất tài liệu chuẩn và sinh diễn giải có kiểm tra. Dự án được phát triển cho cuộc thi DENSO Factory Hacks 2026~\citep{denso}. Phạm vi dữ liệu là nhiệm vụ lắp ráp bộ bánh răng hành tinh của bộ dữ liệu HA4M~\citep{ha4m} với một góc camera cố định; tại thời điểm thực nghiệm, nhóm có ba video có nhãn của hai người (hai video của U001 và một video của U002), nên kết quả được trình bày như một nghiên cứu thí điểm, với người huấn luyện và người kiểm tra tách riêng. Thời gian chuẩn lấy từ tài liệu hướng dẫn công việc do nhóm biên soạn mô phỏng, dùng để kiểm chứng chức năng đối chiếu chứ không đại diện cho định mức của nhà máy.

\subsection{Contributions}

Đóng góp của đề tài mang tính thực nghiệm: (1) một luồng xử lý hoàn chỉnh, có lưu vết từ khung hình đến khuyến nghị; (2) một giao thức đánh giá tách người khi số người rất ít — kiểm định chéo theo video trên người huấn luyện để chọn siêu tham số, kiểm tra trên người mới, 5 seed cho mỗi cấu hình — cùng so sánh với kết quả kiểm tra trong cùng người; (3) phân tích nguyên nhân suy giảm khái quát (găng tay làm hỏng điểm bàn tay, khác biệt tốc độ thao tác) và các biện pháp tương ứng — loại nhóm đặc trưng kém tin cậy, tăng cường co giãn tốc độ, giải mã ràng buộc thứ tự SOP — được đo bằng nghiên cứu loại bỏ thành phần; (4) thiết kế và phân tích định tính tầng RAG, gồm cơ chế kiểm tra đầu ra và đối chiếu phát hiện giữa ranh giới dự đoán và ranh giới chuẩn.
